\documentclass{beamer}
\usepackage[siunitx,europeanresistors,americaninductors]{circuitikz}% elektriskās shēmas
\usepackage{tikz}% priekš elektriskajām shēmām kaut kas
\usepackage{amsmath}% var izmantot math mode
\usepackage{color}% var mainīt krāsu teksta paragrāfa vidū
\usepackage{wrapfig}% attelam apkart iet teksts
\usepackage{blindtext}% labākais veids kā saprast erorrus
\usepackage[utf8]{inputenc}
\usepackage{graphicx} % var ielikt bildes
\usepackage{tabularx} % tabulas


\usepackage{multicol} % vairākas kolonas var izmantot


\usetheme{Goettingen} % Uzliek labjā malā nosaukumu un autoru

\title{Mezglu spriegumu metode (MSM)}
\author{Dans Laksis}
\institute{\LaTeX}
\date{2019}
 
\begin{document} % dokumenta sākums
 
\frame{\titlepage} % titullapa
 
 
%% 1.lekcijas lapa
\input{1.lapa/Oma_likums.tex}
\input{1.lapa/Ievads_MSM.tex}
\input{1.lapa/Piemers_1.tex}


%% 2.lekcijas lapa
\input{2.lapa/Piemers_2.tex}


%% 3.lekcijas lapa
\input{3.lapa/MSM_algoritms.tex}


%@ 4.lekcijas lapa
\input{4.lapa/Aprekina_piemers.tex}



\end{document}